%======================================================================
\subsection{Mates and motion, SolidWorks-style --- Plan for the fork work (2 Nov 2026)}\label{sec:mates-motion-plan}
%======================================================================

You asked, on 2 Nov: \emph{``after I join two parts together I should see them move with each other. Like how
fusion or solidworks moves they keep those constraints. VectorLabs should be more constraint based instead of the
rigid joints''} and \emph{``make sure we touch motion to have those capabilities too''}. Your choices: a pop-up
mate bar, a first scope of Standard + Advanced + Mechanical mates with a degrees-of-freedom (DOF) readout, and the
TL printer's Fixed joints converted automatically.

The work is split three ways. The Vector Labs session builds the \textbf{mate bar} (pick two faces, edges or
points; the mates that fit are offered, with a value, Flip and OK) and the \textbf{printer conversion}. This plan
covers the \textbf{fork work}, B1--B3 below. B4 (the right-panel drawer stays closed while a command is open) and
B5 (a dragged part stays under the cursor) are already built and installed (1.1.4-bw.1+5de7697d).

\paragraph{What the solver already does (read in the source, 2 Nov).}
\begin{itemize}
  \item The standard mates exist as joints. A \emph{Distance} joint picks its solver constraint from the two
        picked shapes (\file{AssemblyObject.cpp}, \code{makeMbdJointDistance}): plane--plane is a planar joint
        (Coincident or Distance), plane--cylinder is Tangent, point--plane is a point on a face, and so on.
        Parallel, Perpendicular, Angle, Fixed (Lock), Revolute (Hinge), Slider, Cylindrical, Ball, Gears, Rack \&
        pinion, Screw and Belt are joints of their own. Several joints between two parts stack.
  \item \textbf{One trap for the mate bar:} a cylinder--cylinder Distance of 0 means the cylinders \emph{touch}
        (the solver adds both radii), not that they are concentric. Concentric is the Cylindrical joint;
        Concentric + Coincident is the Revolute (Hinge). The Vector Labs session has been told.
  \item The solver reports no degrees of freedom. The assembly's ``Under-constrained: $n$ DOF'' text exists, but
        its number is a placeholder (6 per part) that is never computed, and \code{solve()} returns only
        ``solved'', ``nothing grounded'' or ``failed''.
  \item A motion study can drive only Revolute, Slider and Cylindrical joints. Joint limits are left out of the
        solver while a study runs (to avoid clashes with the driven motion), so a run can pass a limit. Gears,
        Rack \& pinion, Screw and Belt are part of the run and carry the motion through.
  \item Nothing in the solver centres one face pair between another (Width) or mirrors two parts about a plane
        (Symmetric).
\end{itemize}

%----------------------------------------------------------------------
\subsubsection*{B2. The DOF readout (first)}

\emph{What you see.} Every component the assembly holds shows \textbf{(-)} in the tree while it can still move,
\textbf{(+)} when a joint on it is redundant or conflicting, and nothing when it is fully defined, as in
SolidWorks. Selecting it, the Properties tab says how it can move: ``Free: 1 rotation, 1 translation'', and
``moves with Part~7'' when gears or a belt tie it to another part. The assembly row says ``Fully defined'' or
``Under-defined: 4 DOF''.

\emph{Theory.} Give every part a small motion: a velocity $\mathbf v_i$ and a spin $\boldsymbol\omega_i$, six
numbers per part, stacked into $\dot{\mathbf q}$. Each joint is a set of equations $\Phi(\mathbf q)=0$; a motion
keeps them true when $\Phi_{\mathbf q}\,\dot{\mathbf q}=0$, where $\Phi_{\mathbf q}$ is the joints' Jacobian.
The motions the assembly still allows are the null space of $\Phi_{\mathbf q}$ (grounded parts fixed). Take a basis
$N$ of that null space and its six rows $N_i$ for part $i$:
\[
  \mathrm{DOF}_i=\operatorname{rank} N_i,\qquad \text{rotations}_i=\operatorname{rank} N_i^{\omega},\qquad
  \text{translations}_i=\mathrm{DOF}_i-\text{rotations}_i ,
\]
where $N_i^{\omega}$ is the spin half of $N_i$. The assembly's DOF is $\operatorname{rank} N$. Two parts share a
column of $N$ when their motions are tied (gears): that gives ``moves with''. Rows of $\Phi_{\mathbf q}$ that are
linear combinations of others are the redundant joints, the (+). The ranks come from a singular value
decomposition with a tolerance scaled to the assembly's size.

\emph{How.} The solver builds $\Phi_{\mathbf q}$ to solve positions but keeps it private. I add one call to our
copy of OndselSolver that returns it at the solved position, compute the ranks in \file{AssemblyObject.cpp} after
each solve (not during a drag), and fill the existing DOF and redundancy fields, so the assembly's message becomes
true. The tree mark uses the Browser's row hook (no tree change); the Properties line is a row in the Properties
tab. First a spike: if the Jacobian cannot be reached cleanly, the fallback probes each free part with the drag
solver (six small trial moves; a move the joints allow is a freedom), which needs no solver change but is slower.

\emph{Cost.} The printer has a few hundred parts, so $\Phi_{\mathbf q}$ has some 2000 columns; the decomposition
takes about a second. It runs after a solve, is skipped while dragging, and is cached until a joint changes.

\emph{Tests.} A box on a Slider (1 translation), on a Revolute (1 rotation), on a Cylindrical (1 + 1), on a
Ball (3 rotations), on a planar Distance (2 translations, 1 rotation), a gear pair (1, moves with), a part with
two Fixed joints (+), and a fully fixed part (no mark).

%----------------------------------------------------------------------
\subsubsection*{B3. Motion on the new mates}

Built on the motion study you already have since 27 Oct (Drive Joint, Motion Study, Motion Link).
\begin{enumerate}
  \item \textbf{Motors on more mates.} Revolute, Slider and Cylindrical keep the solver's kinematic run. A
        Distance or Angle mate gets a \emph{driven value}: the study sets the mate's distance or angle frame by
        frame and solves, so a ``Distance 0~$\rightarrow$~300\,mm'' motor works on any mate the bar makes. Both kinds
        play back through the same timeline.
  \item \textbf{Limits respected.} Limits stay out of the solver during a run (they clash with a driven motion),
        but every frame is checked. The run stops at the first frame past a limit and says which:
        ``Carriage slider reaches its 287\,mm limit at 0.96\,s''. This is what SolidWorks does at a limit.
  \item \textbf{Gears, rack, screw, belt} already carry the motion; tested here with a driven gear pair and a rack.
  \item \textbf{Followers.} Parts mated to a moving part follow it, because they are in the same solve; tested
        with screws on hinge mates riding a carriage.
\end{enumerate}
\emph{Acceptance} (on a test model first; on the printer once its joints are converted): a rail carriage on a Slider
driven 0~$\rightarrow$~300\,mm with its four screws (Hinge mates) riding along, each screw keeping its place on the
carriage; then a gear pair driven by an angular motor, the second gear turning by the ratio.

%----------------------------------------------------------------------
\subsubsection*{B1. Width and Symmetric (last: the most uncertain)}

\begin{figure}[!htbp]
\centering
\begin{tikzpicture}[scale=0.9,>=Stealth]
  % slot (part 1)
  \fill[black!8] (-3,0) rectangle (-1.2,2.4);
  \fill[black!8] (1.2,0) rectangle (3,2.4);
  \draw[thick] (-3,2.4) -- (-1.2,2.4) -- (-1.2,0.4) -- (1.2,0.4) -- (1.2,2.4) -- (3,2.4);
  \node[font=\small] at (-2.1,1.2) {part 1};
  \node[font=\small,bwgrey] at (-1.2,2.75) {$A_1$};
  \node[font=\small,bwgrey] at (1.2,2.75) {$A_2$};
  % tab (part 2)
  \fill[bwblue!15] (-0.6,1.0) rectangle (0.6,3.2);
  \draw[thick,bwblue] (-0.6,1.0) rectangle (0.6,3.2);
  \node[font=\small,bwblue,anchor=east] at (-0.7,3.0) {part 2};
  \node[font=\small,bwblue] at (-0.85,1.55) {$B_1$};
  \node[font=\small,bwblue] at (0.85,1.55) {$B_2$};
  % midplanes
  \draw[dashed,bworange,thick] (0,-0.3) -- (0,3.9);
  \node[font=\small,bworange] at (0,-0.65) {midplane of $A_1A_2$ = midplane of $B_1B_2$};
  \draw[<->] (-1.2,0.15) -- (0,0.15);
  \draw[<->] (0,0.15) -- (1.2,0.15);
\end{tikzpicture}
\caption{Width: the tab's two faces $B_1,B_2$ are centred between the slot's faces $A_1,A_2$. Each face pair is on
one rigid part, so its midplane is a fixed frame on that part; the mate is a planar joint between the two midplanes.}
\label{fig:width-mate}
\end{figure}

\textbf{Width} (figure~\ref{fig:width-mate}) needs no solver change. The mate keeps two faces on each side; the
joint's two frames are the two midplanes, computed where every joint computes its frames today
(\file{JointObject.py}), and the solver gets a planar joint between them. The tab can still slide along the slot
and lift out of it, as in SolidWorks' centred Width. Its new joint type is added at the end of the list (stored
files keep their numbers), with an icon, the mate bar entry and a test: a tab centred in slots of 20, 30 and 40\,mm
wide stays centred while it is dragged.

\textbf{Symmetric} mirrors an entity on one part and an entity on another about a plane.
\begin{itemize}
  \item When the plane is on one of the two parts, the mirror image is a fixed frame on that part, and the mate is
        an ordinary joint between the other part's entity and that image. No solver change; built with Width.
  \item The usual case, though, is three parts: two vise jaws mirrored about the centre plane of the base. Both
        jaws move, so their motions are tied (one moves out by as much as the other). The solver's joints join
        two parts each, and its couplings (gears, rack, screw) tie a rotation to a rotation or to a translation,
        never two translations. This needs a new solver constraint: point--point and plane--plane symmetric about
        a plane on a third part, three equations each. That is real solver work, about two to three days, with the
        Jacobian derivatives tested against finite differences.
\end{itemize}

%----------------------------------------------------------------------
\subsubsection*{Order and estimate}
\begin{center}\small
\begin{tabularx}{\linewidth}{@{}lXl@{}}
\toprule
Step & What & Estimate\\
\midrule
B2 & Spike (reach the Jacobian), then the DOF readout: tree (-)/(+), Properties line, assembly message & half a day to a day\\
B3 & Driven Distance and Angle mates, limits checked every frame, gear/rack/screw/follower tests, acceptance model & a day\\
B1a & Width, and Symmetric with the plane on one of the two parts & a day\\
B1b & Symmetric across three parts (new solver constraint), if you want it now & two to three days\\
\bottomrule
\end{tabularx}
\end{center}
Each step is committed and staged when it passes its tests, as usual; nothing waits for the whole set.

\paragraph{Questions for you.}
\begin{enumerate}
  \item The \textbf{(-)} mark: before the name, as in SolidWorks (``(-) TL-01-01147 Carriage''), which needs a
        small change to how the tree draws a row; or as a small chip after the name, which the tree can already
        draw? I recommend before the name.
  \item \textbf{Symmetric across three parts} (the vise case): build it now (B1b, two to three days of solver work),
        or first Width and the two-part Symmetric (B1a) and the three-part case later? I recommend B1a first.
  \item \textbf{A motion that reaches a limit}: stop the run there and say which limit (recommended, as
        SolidWorks), or hold that joint at its limit and let the rest of the run continue?
  \item \textbf{Width options}: centred only for now (recommended), or also SolidWorks' ``free'' (the tab may move
        between the walls), ``distance from one wall'' and ``percent'' variants?
  \item \textbf{DOF inside sub-assemblies}: marks on the top assembly's own components and on the parts of
        flexible sub-assemblies (recommended), or on every level, rigid sub-assemblies included (they cannot move
        inside, so they would all read fully defined)?
\end{enumerate}
